\section{建模行为与强化学习基础}

\subsection{经验的数据表示}

\begin{importantbox}{核心术语}
\begin{itemize}
    \item \textbf{状态} $s_t$：时刻 $t$ 的世界状态。
    \item \textbf{观察} $o_t$：智能体在时刻 $t$ 观察到的信息（当存在信息缺失时使用）。
    \item \textbf{动作} $a_t$：时刻 $t$ 的决策。
    \item \textbf{轨迹} $\tau = (s_1, a_1, s_2, a_2, \ldots, s_T, a_T)$：一个完整的状态-动作序列。
    \item \textbf{奖励函数} $r(s, a)$：衡量状态-动作对的好坏。
\end{itemize}
\end{importantbox}

\subsection{状态与观察的区别}

一个关键性质是\textbf{马尔可夫性}（Markov Property）：下一个状态仅取决于当前状态和动作，与更早的历史无关：
$$p(s_{t+1} \mid s_t, a_t) \quad \text{独立于 } s_{t-1}$$

但在很多实际场景中（如聊天机器人），智能体只能看到\textbf{观察}而非完整状态，此时需要给策略加入记忆：$\pi_\theta(a_t \mid o_{t-m}, \ldots, o_t)$。

\subsection{用神经网络表示行为}

策略 $\pi_\theta(a \mid s)$ 是将状态映射到动作分布的函数。智能体的运行过程如下：
\begin{enumerate}
    \item 观察状态 $s_t$
    \item 从策略中采样动作 $a_t \sim \pi_\theta(\cdot \mid s_t)$
    \item 观察下一个状态 $s_{t+1} \sim p(\cdot \mid s_t, a_t)$
\end{enumerate}
这样产生的完整序列称为策略 rollout 或 episode。

\subsection{强化学习的目标}

强化学习的目标是\textbf{最大化期望累计奖励}：
$$\max_\theta \; \mathbb{E}_{\tau \sim p_\theta(\tau)} \left[ \sum_{t} r(s_t, a_t) \right]$$

其中轨迹分布为：
$$p_\theta(\tau) = p(s_1) \prod_{t=1}^{T} \pi_\theta(a_t \mid s_t) \, p(s_{t+1} \mid s_t, a_t)$$

\begin{warningbox}{为什么要用期望}
累计奖励不是一个确定量，它有两个随机性来源：(1) 环境本身是随机的；(2) 策略可能是随机的。因此我们优化的是\textbf{期望}累计奖励。
\end{warningbox}

\subsection{随机策略的意义}

为什么使用随机策略而不是确定性策略？两个原因：
\begin{enumerate}
    \item \textbf{探索}：要从自己的经验中学习，必须尝试不同的事情。
    \item \textbf{建模随机行为}：现有数据往往展示多样化的行为，可以利用生成模型的工具。
\end{enumerate}

\subsection{价值函数与 Q 函数}

衡量策略好坏的两个核心量：

\begin{itemize}
    \item \textbf{价值函数} $V^\pi(s)$：从状态 $s$ 出发，遵循策略 $\pi$ 的未来期望奖励。
    \item \textbf{Q 函数} $Q^\pi(s, a)$：从状态 $s$ 采取动作 $a$，然后遵循策略 $\pi$ 的未来期望奖励。
\end{itemize}

\subsection{算法类型概览}

\begin{enumerate}
    \item \textbf{Imitation Learning}：模仿能获得高奖励的策略。
    \item \textbf{Policy Gradients}：直接对目标函数求梯度。
    \item \textbf{Actor-Critic}：估计当前策略的价值，并用它来改进策略。
    \item \textbf{Value-based}：估计最优策略的价值。
    \item \textbf{Model-based}：学习环境动力学模型，用于规划或策略改进。
\end{enumerate}

不同算法做了不同的权衡，适用于不同的假设条件（数据收集成本、监督形式、稳定性、动作空间特性等）。

\subsection{本章小结}

强化学习的基本要素包括状态、动作、轨迹、奖励和策略。目标是最大化期望累计奖励。价值函数和 Q 函数是评估策略好坏的核心工具。不同类型的算法各有侧重，适用于不同场景。

